% Technical Documentation Appendix - Code Samples

\section{AWS Serverless Architecture}

\subsection{Full-Stack Serverless Application Overview}

The flagship cloud project demonstrates a complete serverless architecture integrating React frontend, API Gateway, Lambda, and DynamoDB.

\subsubsection{Architecture Diagram}

\begin{lstlisting}[caption=Serverless Application Architecture, captionpos=b]
CLIENT LAYER
=====================================
React Static Website (Hosted on S3 + CloudFront)
- User Interface for displaying records
- Handles GET requests to API Gateway
=====================================
                    |
                    | HTTPS Request
                    v
API LAYER
=====================================
API Gateway REST Endpoint (GET /records)
- Routes incoming HTTP requests
- Handles CORS and request validation
- Authentication and authorization
=====================================
                    |
                    | Invoke
                    v
COMPUTE LAYER (Serverless)
=====================================
AWS Lambda Function (Node.js Runtime)
- Processes API Gateway events
- Queries DynamoDB for records
- Formats and returns response
- Auto-scales based on demand
=====================================
                    |
                    | Query
                    v
DATA LAYER
=====================================
DynamoDB Table (Serverless NoSQL Database)
- Stores application records
- Auto-scaling throughput
- High availability across AZs
- On-demand billing model
=====================================
\end{lstlisting}

\subsubsection{Technology Stack Summary}

\begin{table}[h]
    \centering
    \begin{tabularx}{\textwidth}{|l|l|X|}
        \hline
        \textbf{Layer} & \textbf{Technology} & \textbf{Purpose} \\
        \hline
        Presentation & React + S3 + CloudFront & Serve static website with CDN caching \\
        \hline
        API & API Gateway & HTTP endpoint routing and request handling \\
        \hline
        Logic & AWS Lambda (Node.js) & Serverless function execution \\
        \hline
        Data & DynamoDB & NoSQL database with auto-scaling \\
        \hline
    \end{tabularx}
    \caption{Serverless Architecture Technology Stack}
\end{table}

\subsection{Lambda Function Implementation}

\begin{lstlisting}[basicstyle=\ttfamily\small, caption=AWS Lambda Handler for GET Records, captionpos=b]
// Lambda handler function for API Gateway GET request
const AWS = require('aws-sdk');
const dynamoDB = new AWS.DynamoDB.DocumentClient();

const TABLE_NAME = process.env.DYNAMODB_TABLE;

exports.handler = async (event) => {
  try {
    console.log('Received event:', JSON.stringify(event));
    
    // Parse query parameters
    const { recordId, limit = 10 } = event.queryStringParameters || {};
    
    // Query DynamoDB based on parameters
    let params;
    if (recordId) {
      // Get single record by ID
      params = {
        TableName: TABLE_NAME,
        Key: { id: recordId }
      };
      const result = await dynamoDB.get(params).promise();
      return formatResponse(200, result.Item);
    } else {
      // Scan all records with limit
      params = {
        TableName: TABLE_NAME,
        Limit: limit
      };
      const result = await dynamoDB.scan(params).promise();
      return formatResponse(200, result.Items);
    }
  } catch (error) {
    console.error('Error:', error);
    return formatResponse(500, { 
      error: 'Internal server error',
      message: error.message 
    });
  }
};

// Helper function to format API responses
const formatResponse = (statusCode, body) => {
  return {
    statusCode,
    headers: {
      'Content-Type': 'application/json',
      'Access-Control-Allow-Origin': '*'
    },
    body: JSON.stringify(body)
  };
};
\end{lstlisting}

\subsection{DynamoDB Table Configuration}

\begin{lstlisting}[basicstyle=\ttfamily\small, caption=DynamoDB Table Schema Setup, captionpos=b]
// DynamoDB table creation with CloudFormation or AWS SDK
const tableParams = {
  TableName: 'ApplicationRecords',
  KeySchema: [
    { AttributeName: 'id', KeyType: 'HASH' }  // Partition key
  ],
  AttributeDefinitions: [
    { AttributeName: 'id', AttributeType: 'S' },
    { AttributeName: 'timestamp', AttributeType: 'N' }
  ],
  BillingMode: 'PAY_PER_REQUEST',  // Auto-scaling
  GlobalSecondaryIndexes: [
    {
      IndexName: 'TimestampIndex',
      KeySchema: [
        { AttributeName: 'timestamp', KeyType: 'HASH' }
      ],
      Projection: { ProjectionType: 'ALL' },
      BillingMode: 'PAY_PER_REQUEST'
    }
  ],
  Tags: [
    { Key: 'Application', Value: 'ServerlessApp' },
    { Key: 'Environment', Value: 'Production' }
  ]
};

// Sample record structure
const sampleRecord = {
  id: 'record-001',
  title: 'Sample Record',
  description: 'This is a sample DynamoDB record',
  timestamp: Date.now(),
  status: 'active',
  metadata: {
    createdBy: 'system',
    lastModified: Date.now()
  }
};
\end{lstlisting}

\subsection{React Frontend Integration}

\begin{lstlisting}[basicstyle=\ttfamily\small, caption=React Component Calling Lambda via API Gateway, captionpos=b]
// React component fetching records from Lambda via API Gateway
import React, { useState, useEffect } from 'react';

function RecordsList() {
  const [records, setRecords] = useState([]);
  const [loading, setLoading] = useState(true);
  const [error, setError] = useState(null);

  useEffect(() => {
    const fetchRecords = async () => {
      try {
        const apiEndpoint = process.env.REACT_APP_API_ENDPOINT;
        const response = await fetch(
          `${apiEndpoint}/records?limit=20`,
          {
            method: 'GET',
            headers: {
              'Content-Type': 'application/json'
            }
          }
        );

        if (!response.ok) {
          throw new Error(`HTTP error! status: ${response.status}`);
        }

        const data = await response.json();
        setRecords(data);
        setLoading(false);
      } catch (err) {
        setError(err.message);
        setLoading(false);
      }
    };

    fetchRecords();
  }, []);

  if (loading) return <div>Loading records...</div>;
  if (error) return <div>Error: {error}</div>;

  return (
    <div className="records-container">
      <h1>Application Records</h1>
      <ul>
        {records.map((record) => (
          <li key={record.id}>
            <h3>{record.title}</h3>
            <p>{record.description}</p>
            <small>Created: {new Date(record.timestamp).toLocaleDateString()}</small>
          </li>
        ))}
      </ul>
    </div>
  );
}

export default RecordsList;
\end{lstlisting}

\section{Key Implementation Examples}

\subsection{JWT Authentication Middleware}

\begin{lstlisting}[basicstyle=\ttfamily\small, caption=Express JWT Middleware, captionpos=b]
// Authentication middleware for protected routes
const verifyToken = (req, res, next) => {
  const token = req.headers['authorization']?.split(' ')[1];
  
  if (!token) return res.status(403).json({ 
    message: 'No token provided' 
  });
  
  jwt.verify(token, process.env.JWT_SECRET, (err, decoded) => {
    if (err) return res.status(401).json({ 
      message: 'Invalid token' 
    });
    
    req.userId = decoded.id;
    next();
  });
};

// Usage: app.get('/protected', verifyToken, (req, res) => {...})
\end{lstlisting}

\subsection{Docker Compose Configuration}

\begin{lstlisting}[language=yaml, caption=Multi-Container Docker Setup, captionpos=b]
version: '3.8'
services:
  app:
    build: .
    ports:
      - "3000:3000"
    environment:
      - MONGODB_URI=mongodb://mongo:27017/appdb
    depends_on:
      - mongo
    networks:
      - app-network
  
  mongo:
    image: mongo:5.0
    environment:
      MONGO_INITDB_ROOT_USERNAME: admin
      MONGO_INITDB_ROOT_PASSWORD: password
    volumes:
      - mongo_data:/data/db
    networks:
      - app-network
  
  mongo-express:
    image: mongo-express
    ports:
      - "8081:8081"
    environment:
      ME_CONFIG_MONGODB_ADMINUSERNAME: admin
      ME_CONFIG_MONGODB_ADMINPASSWORD: password
    depends_on:
      - mongo
    networks:
      - app-network

volumes:
  mongo_data:

networks:
  app-network:
    driver: bridge
\end{lstlisting}

\subsection{Real-Time Chat Event Handling}

\begin{lstlisting}[basicstyle=\ttfamily\small, caption=Socket.IO Event Listeners, captionpos=b]
io.on('connection', (socket) => {
  console.log(`User ${socket.id} connected`);
  
  // User joins a room
  socket.on('join-room', ({ room, username }) => {
    socket.join(room);
    socket.username = username;
    socket.room = room;
    
    io.to(room).emit('user-joined', {
      username,
      message: `${username} joined the room`
    });
  });
  
  // User sends a message
  socket.on('send-message', ({ message }) => {
    io.to(socket.room).emit('receive-message', {
      username: socket.username,
      message,
      timestamp: new Date()
    });
  });
  
  // User leaves room
  socket.on('leave-room', () => {
    io.to(socket.room).emit('user-left', {
      username: socket.username,
      message: `${socket.username} left the room`
    });
    socket.leave(socket.room);
  });
  
  socket.on('disconnect', () => {
    console.log(`User ${socket.id} disconnected`);
  });
});
\end{lstlisting}

\subsection{Redis Caching Implementation}

\begin{lstlisting}[basicstyle=\ttfamily\small, caption=Redis Cache with TTL, captionpos=b]
const redis = require('redis');
const client = redis.createClient({
  host: 'localhost',
  port: 6379
});

client.connect();

const getCachedData = async (key) => {
  try {
    // Try to get from cache
    const cached = await client.get(key);
    if (cached) {
      console.log('Cache HIT');
      return JSON.parse(cached);
    }
    
    // Cache miss - fetch from API
    console.log('Cache MISS');
    const data = await fetch('https://api.example.com/data').then(
      r => r.json()
    );
    
    // Store in cache with 3600 second TTL
    await client.setEx(key, 3600, JSON.stringify(data));
    return data;
  } catch (error) {
    console.error('Cache error:', error);
    throw error;
  }
};
\end{lstlisting}

\subsection{Kafka Producer-Consumer Pattern}

\begin{lstlisting}[basicstyle=\ttfamily\small, caption=Kafka Message Handling, captionpos=b]
const { Kafka } = require('kafkajs');

const kafka = new Kafka({
  clientId: 'my-app',
  brokers: ['localhost:9092']
});

// Producer
const producer = kafka.producer();

const publishEvent = async (topic, message) => {
  await producer.connect();
  await producer.send({
    topic,
    messages: [
      { key: message.id, value: JSON.stringify(message) }
    ],
  });
  await producer.disconnect();
};

// Consumer
const consumer = kafka.consumer({ groupId: 'my-group' });

const subscribeToEvents = async (topic, handler) => {
  await consumer.connect();
  await consumer.subscribe({ topic, fromBeginning: true });
  
  await consumer.run({
    eachMessage: async ({ topic, partition, message }) => {
      const data = JSON.parse(message.value.toString());
      await handler(data);
    },
  });
};
\end{lstlisting}

\subsection{MongoDB Transaction Example}

\begin{lstlisting}[basicstyle=\ttfamily\small, caption=ACID Transaction, captionpos=b]
const performTransfer = async (fromUserId, toUserId, amount) => {
  const session = await db.startSession();
  try {
    await session.withTransaction(async () => {
      // Deduct from sender
      await User.updateOne(
        { _id: fromUserId },
        { $inc: { balance: -amount } },
        { session }
      );
      
      // Add to receiver
      await User.updateOne(
        { _id: toUserId },
        { $inc: { balance: amount } },
        { session }
      );
      
      // Create transaction record
      await Transaction.create([{
        from: fromUserId,
        to: toUserId,
        amount,
        timestamp: new Date()
      }], { session });
    });
  } finally {
    await session.endSession();
  }
};
\end{lstlisting}

\section{Project Structure Overview}

\subsection{Recommended Directory Layout}

\begin{lstlisting}[language=bash, caption=Backend Project Structure, captionpos=b]
project-root/
  ├── src/
  │   ├── config/           # Configuration files
  │   ├── controllers/      # Route handlers
  │   ├── middlewares/      # Express middlewares
  │   ├── models/           # Database models
  │   ├── routes/           # Route definitions
  │   ├── services/         # Business logic
  │   ├── utils/            # Helper functions
  │   └── app.js            # Express app setup
  ├── tests/                # Unit and integration tests
  ├── .env                  # Environment variables
  ├── .gitignore            # Git ignore file
  ├── docker-compose.yml    # Container orchestration
  ├── Dockerfile            # Docker image definition
  ├── package.json          # Dependencies
  └── README.md             # Project documentation
\end{lstlisting}

\subsection{Serverless Project Structure}

\begin{lstlisting}[language=bash, caption=AWS Serverless Application Structure, captionpos=b]
serverless-app/
  ├── frontend/             # React application
  │   ├── src/
  │   │   ├── components/
  │   │   ├── pages/
  │   │   └── App.js
  │   ├── public/
  │   └── package.json
  │
  ├── backend/              # Lambda functions
  │   ├── functions/
  │   │   ├── getRecords.js     # GET handler
  │   │   ├── createRecord.js   # POST handler
  │   │   └── utils/
  │   └── package.json
  │
  ├── .github/
  │   └── workflows/
  │       └── deploy.yml    # GitHub Actions CI/CD
  │
  ├── .ebextensions/        # Elastic Beanstalk config
  │   └── deployer.config
  │
  ├── serverless.yml        # Serverless Framework config
  └── README.md
\end{lstlisting}

\section{CI/CD Automation with Elastic Beanstalk}

\subsection{GitHub Actions Workflow}

\begin{lstlisting}[language=yaml, caption=GitHub Actions CI/CD Pipeline, captionpos=b]
name: Deploy to AWS

on:
  push:
    branches: [ main, develop ]
  pull_request:
    branches: [ main ]

jobs:
  build-and-deploy:
    runs-on: ubuntu-latest

    steps:
    - uses: actions/checkout@v2

    - name: Setup Node.js
      uses: actions/setup-node@v2
      with:
        node-version: '16'

    # Build React frontend
    - name: Build React Frontend
      working-directory: ./frontend
      run: |
        npm install
        npm run build
        echo "Frontend built successfully"

    # Test Lambda functions
    - name: Test Lambda Functions
      working-directory: ./backend
      run: |
        npm install
        npm test
        echo "Tests passed"

    # Package Lambda function
    - name: Package Lambda Function
      working-directory: ./backend
      run: |
        zip -r lambda-function.zip .
        echo "Lambda packaged"

    # Deploy to AWS Elastic Beanstalk
    - name: Deploy to Elastic Beanstalk
      if: github.ref == 'refs/heads/main'
      env:
        AWS_ACCESS_KEY_ID: ${{ secrets.AWS_ACCESS_KEY_ID }}
        AWS_SECRET_ACCESS_KEY: ${{ secrets.AWS_SECRET_ACCESS_KEY }}
      run: |
        aws elasticbeanstalk create-application-version \
          --application-name serverless-app \
          --version-label ${{ github.sha }} \
          --source-bundle S3Bucket=${{ secrets.EB_BUCKET }},S3Key=frontend-build.zip

        aws elasticbeanstalk update-environment \
          --environment-name production \
          --version-label ${{ github.sha }}

        echo "Deployment initiated"

    - name: Health Check
      run: |
        sleep 30
        STATUS=$(curl -s -o /dev/null -w "%{http_code}" ${{ secrets.EB_URL }})
        if [ "$STATUS" = "200" ]; then
          echo "Deployment successful"
        else
          echo "Deployment health check failed"
          exit 1
        fi
\end{lstlisting}

\subsection{Elastic Beanstalk Configuration}

\begin{lstlisting}[language=yaml, caption=AWS Elastic Beanstalk Configuration, captionpos=b]
# .ebextensions/deployer.config
option_settings:
  aws:autoscaling:asg:
    MaxSize: 6
    MinSize: 2
  aws:autoscaling:launchconfiguration:
    InstanceType: t3.micro
  aws:ec2:instances:
    InstanceTypes: t3.micro
  aws:elasticbeanstalk:application:
    Application Healthcheck URL: /health
  aws:elasticbeanstalk:environment:
    EnvironmentType: LoadBalanced
    LoadBalancerType: application
  aws:elasticbeanstalk:cloudwatch:logs:
    StreamLogs: true
    DeleteOnTerminate: false
    RetentionInDays: 7
  aws:elasticbeanstalk:xray:
    XRayEnabled: true

commands:
  01_install_node_dependencies:
    command: "npm install"
    leader_only: true
  02_build_react:
    command: "npm run build"
    leader_only: true

container_commands:
  01_migrate:
    command: "npm run migrate"
    leader_only: true
  02_seed_db:
    command: "npm run seed"
    leader_only: true
\end{lstlisting}

\subsection{Deployment Process Flow}

\begin{enumerate}
    \item \textbf{Code Commit:} Developer pushes to GitHub main branch
    \item \textbf{Trigger:} GitHub Actions workflow automatically triggered
    \item \textbf{Build Stage:} 
    \begin{itemize}
        \item Install dependencies
        \item Build React frontend (optimized for production)
        \item Run unit tests for Lambda
    \end{itemize}
    \item \textbf{Package Stage:}
    \begin{itemize}
        \item Create deployment artifact
        \item Upload to S3 bucket
    \end{itemize}
    \item \textbf{Deploy Stage:}
    \begin{itemize}
        \item Create new application version in Elastic Beanstalk
        \item Update environment to new version
        \item Perform rolling deployment
    \end{itemize}
    \item \textbf{Verification:}
    \begin{itemize}
        \item Health checks verify deployment success
        \item Automatic rollback on failure
    \end{itemize}
\end{enumerate}

\subsection{Environment Variables for CI/CD}

\begin{lstlisting}[language=bash, caption=Required GitHub Secrets for Deployment, captionpos=b]
# GitHub Secrets to configure:
AWS_ACCESS_KEY_ID              # AWS IAM user access key
AWS_SECRET_ACCESS_KEY          # AWS IAM user secret key
AWS_REGION                     # e.g., us-east-1
EB_BUCKET                      # S3 bucket for Elastic Beanstalk artifacts
EB_APPLICATION_NAME            # Elastic Beanstalk app name
EB_ENVIRONMENT_NAME            # Environment (dev/staging/prod)
EB_URL                         # Environment URL for health checks
DYNAMODB_TABLE_NAME            # DynamoDB table name
API_ENDPOINT                   # API Gateway endpoint URL

# These are injected into the deployment pipeline automatically
\end{lstlisting}

\section{Performance Benchmark Results}

\subsection{Serverless Architecture Performance}

\begin{table}[h]
    \centering
    \begin{tabular}{|l|c|c|}
        \hline
        \textbf{Metric} & \textbf{Value} & \textbf{Implication} \\
        \hline
        Lambda Cold Start & 500-800ms & First execution after inactivity \\
        \hline
        Lambda Warm Start & 10-50ms & Subsequent executions \\
        \hline
        API Gateway Latency & 50-100ms & Request routing overhead \\
        \hline
        DynamoDB Query & 10-20ms & Direct table query \\
        \hline
        End-to-End Response & 100-150ms & Full request cycle (warm) \\
        \hline
        Concurrent Executions & 1000+ & Default AWS account limit \\
        \hline
    \end{tabular}
    \caption{AWS Serverless Performance Metrics}
\end{table}

\subsection{Cost Analysis}

\begin{table}[h]
    \centering
    \begin{tabularx}{\textwidth}{|l|X|c|}
        \hline
        \textbf{Service} & \textbf{Pricing Model} & \textbf{Monthly Cost (1M req)} \\
        \hline
        Lambda & \$0.20 per 1M requests + compute time & \$2-5 \\
        \hline
        API Gateway & \$3.50 per 1M API calls & \$3.50 \\
        \hline
        DynamoDB (on-demand) & \$1.25 per 1M write + \$0.25 per 1M read & \$1-3 \\
        \hline
        S3 (frontend) & \$0.023 per GB stored & \$1 \\
        \hline
        CloudFront & \$0.085 per GB delivered & \$2-5 \\
        \hline
        \textbf{TOTAL} & & \textbf{\$10-20} \\
        \hline
    \end{tabularx}
    \caption{Serverless Architecture Monthly Cost Estimate}
\end{table}

\subsection{Redis Caching Effectiveness}

\begin{table}[h]
    \centering
    \begin{tabular}{|l|c|c|c|}
        \hline
        \textbf{Scenario} & \textbf{Without Cache} & \textbf{With Cache} & \textbf{Improvement} \\
        \hline
        Response Time & 250ms & 8ms & 96.8\% faster \\
        \hline
        Server CPU Usage & 85\% & 25\% & 70.6\% reduction \\
        \hline
        Memory Usage & 512MB & 480MB & 6.2\% saving \\
        \hline
        Requests/sec & 500 & 5000 & 10x throughput \\
        \hline
    \end{tabular}
    \caption{Redis Caching Performance Metrics}
\end{table}

\subsection{Docker Deployment Efficiency}

\begin{table}[h]
    \centering
    \begin{tabular}{|l|c|c|c|}
        \hline
        \textbf{Metric} & \textbf{Traditional VM} & \textbf{Docker Container} & \textbf{Improvement} \\
        \hline
        Deployment Time & 45 min & 5 min & 89\% faster \\
        \hline
        Boot Time & 3-5 min & 2-3 sec & 99\% faster \\
        \hline
        Disk Space & 4-8 GB & 100-500 MB & 90\% reduction \\
        \hline
        Resource Overhead & 20-30\% & 5-10\% & 67\% saving \\
        \hline
    \end{tabular}
    \caption{Docker vs Traditional Deployment}
\end{table}

\section{Learning Resources}

\subsection{Recommended Tools and Platforms}

\begin{itemize}
    \item \textbf{Backend Frameworks:} Express.js, Fastify, NestJS
    \item \textbf{Real-Time Libraries:} Socket.IO, WebSocket, GraphQL Subscriptions
    \item \textbf{Message Brokers:} Kafka, RabbitMQ, AWS SQS
    \item \textbf{Cloud Platforms:} AWS, Google Cloud, Azure
    \item \textbf{Container Tools:} Docker, Docker Compose, Kubernetes
    \item \textbf{Monitoring:} Prometheus, Grafana, ELK Stack
\end{itemize}

\subsection{GitHub Repositories for Reference}

All internship project code is available in the main repository with proper documentation:

\begin{itemize}
    \item \texttt{/DAY\_1\&2/} - JavaScript fundamentals
    \item \texttt{/DAY\_3/} - Authentication systems
    \item \texttt{/DAY\_5-6/} - Docker projects
    \item \texttt{/DAY\_8-9/} - Real-time chat applications
    \item \texttt{/DAY\_10/} - Caching and rate limiting
    \item \texttt{/DAY\_13/} - Kafka implementation
    \item \texttt{/DAY\_27-31/} - AWS cloud projects
\end{itemize}
